\appendix

\clearpage

\section{ML Model Details}
\label{app:ml_details}

\subsection{XGBoost}
\label{app:xgboost}

We use the XGBoost gradient-boosted tree implementation with the following default hyperparameters, subject to re-tuning every 24 months:

\begin{table}[h!]
\centering
\begin{tabular}{l l l}
\midrule
Parameter & Default & Grid \\
\midrule
\texttt{max\_depth} & 4 & \{3, 4, 5, 6\} \\
\texttt{learning\_rate} & 0.05 & \{0.01, 0.05, 0.10\} \\
\texttt{n\_estimators} & 300 & \{100, 200, 300, 500\} \\
\texttt{subsample} & 0.8 & \{0.6, 0.8, 1.0\} \\
\texttt{colsample\_bytree} & 0.8 & \{0.6, 0.8, 1.0\} \\
\texttt{reg\_alpha} & 0.1 & \{0.0, 0.1, 1.0\} \\
\texttt{reg\_lambda} & 1.0 & \{0.1, 1.0, 10.0\} \\
\midrule
\end{tabular}
\caption{\footnotesize \bf XGBoost hyperparameters\normalfont. Default values and grid search ranges. Re-tuning occurs every 24 months using the most recent 12-month validation window.}
\label{tab:xgb_params}
\end{table}

Sample weights enter through XGBoost's native gradient weighting mechanism: each observation's contribution to the gradient is scaled by $w_{i,t}$, so that liquid stocks have larger gradients and hence greater influence on the fitted trees.

Early stopping is disabled. In preliminary experiments, we found that early stopping fires too aggressively on noisy return targets, often terminating at iteration 0 or producing near-constant predictions. This is expected behavior: MSE optimisation in low signal-to-noise environments rewards the unconditional mean, and early stopping amplifies this tendency by halting before the model has found cross-sectional patterns.

\subsection{Random Forest}
\label{app:rf}

We use the scikit-learn \texttt{RandomForestRegressor} with:

\begin{table}[h!]
\centering
\begin{tabular}{l l l}
\midrule
Parameter & Default & Grid \\
\midrule
\texttt{n\_estimators} & 300 & \{100, 200, 300, 500\} \\
\texttt{max\_depth} & 6 & \{4, 6, 8, None\} \\
\texttt{max\_features} & 0.33 & \{0.20, 0.33, 0.50\} \\
\texttt{min\_samples\_leaf} & 50 & \{20, 50, 100\} \\
\midrule
\end{tabular}
\caption{\footnotesize \bf Random Forest hyperparameters\normalfont.}
\label{tab:rf_params}
\end{table}

Sample weights enter through the impurity criterion: the variance reduction at each split is computed using weighted observations, so that liquid stocks have greater influence on the split decisions within each tree.

\subsection{Neural Network}
\label{app:nn}

The neural network architecture consists of three hidden layers with 64, 32, and 16 neurons respectively, BatchNormalization after each hidden layer, ReLU activations, and a linear output layer. Training uses the Adam optimiser with:

\begin{itemize}
\item Initial learning rate: 0.001 with cosine annealing
\item Batch size: 4096
\item Maximum epochs: 100 with early stopping (patience=10) on validation loss
\item Dropout: 0.10 after each hidden layer
\end{itemize}

Sample weights enter through Keras's native \texttt{sample\_weight} parameter in the \texttt{fit()} method, which scales the per-observation MSE loss.

SHAP values are computed using DeepExplainer when available; when DeepExplainer fails (e.g., due to incompatible layer types), we fall back to KernelExplainer on a subsample of 500 background observations.


\section{Feature List}
\label{app:features}

Table~\ref{tab:features} lists the 75 stock-level predictors from \citet{chen2022open}, grouped by economic category. We further classify features into ``tradable'' and ``illiquidity'' groups for Hypothesis~\ref{hyp:h2} testing.

% \input{TablesNew/FeatureList}

\paragraph{Illiquidity features (8):} IdioVol3F, IdioVolAHT, zerotrade1M, zerotrade6M, zerotrade12M, MaxRet, VolSD, BetaLiquidityPS.

\paragraph{Tradable features (8):} Mom12m, Mom6m, BMdec, GP, AssetGrowth, RoE, CF, CBOperProf.


\section{Ledoit-Wolf Bootstrap}
\label{app:bootstrap}

We test Sharpe ratio differences using the studentised circular-block bootstrap of \citet{ledoit2008robust}. The procedure is as follows:

\begin{enumerate}
\item Compute the observed Sharpe ratio difference $\hat{\Delta} = \widehat{\SR}_A - \widehat{\SR}_B$ and its HAC standard error $\hat{\sigma}_{\Delta}$ using a prewhitened Parzen kernel.
\item Compute the studentised statistic $d = \hat{\Delta} / \hat{\sigma}_{\Delta}$.
\item Generate $B = 5{,}000$ bootstrap samples using circular blocks of length $b$ (calibrated via Algorithm 3.1 of Ledoit and Wolf, 2008).
\item For each bootstrap sample $b$, compute $d^*_b = \hat{\Delta}^*_b / \hat{\sigma}^*_{\Delta,b}$.
\item The one-sided $p$-value is $p = \frac{1}{B}\sum_{b=1}^{B} \Ind\{d^*_b \geq d\}$.
\end{enumerate}


\section{Additional Empirical Results}
\label{app:additional_results}

% Additional tables and figures not included in the main text.
% \subsection{Factor Alpha Regressions}
% \input{TablesNew/FactorAlphas_Full}

% \subsection{OOS $R^2$ by Model and Training}
% \input{TablesNew/OOSR2}

% \subsection{SHAP Importance Over Time}
% Figure: importance_over_time_{model}.png

% \subsection{Cumulative Returns Under Alternative AUM}
% Figure: cumulative_returns_net_{model}.png
